\documentclass[runningheads]{llncs}
\usepackage[T1]{fontenc}
\usepackage{graphicx}
\usepackage{hyperref}
\usepackage{booktabs}
\usepackage{amsmath}

\begin{document}

\title{Robust Open-Set Plant Disease Recognition using Vision-Language Models}

\author{Anonymous Author(s)}

\institute{Anonymous Institute}

\maketitle

\begin{abstract}
Supervised closed-set architectures perform well on curated benchmarks but struggle in real-world agricultural deployments. Two core problems arise: domain gaps between training images and field conditions, and the inability of fixed-label models to reject diseases or pests not seen during training. This paper presents a multimodal open-set pipeline that separates localization from identification --- Grounding DINO handles open-vocabulary detection of leaves and pests, while a fine-tuned CLIP model classifies each crop region and flags low-confidence predictions as unknown. This design avoids the scale-mismatch failures common in tiling-based approaches such as SAHI, and removes the overconfident misclassifications that closed-set softmax models inevitably produce on out-of-distribution inputs. Evaluated on PlantVillage across 10 known Tomato disease categories and held-out unknown classes, the pipeline achieves an Unknown F1 of 0.526, compared to 0.294 for a ResNet-50 baseline. A few-shot extension allows new disease prototypes to be registered at runtime without retraining, making the system practical for field conditions where new diseases can appear at any time.

\keywords{Plant Disease Detection \and Open-Set Recognition \and Vision-Language Models \and Grounding DINO \and CLIP}
\end{abstract}

\section{Introduction}
Plant disease detection has advanced rapidly on curated benchmarks such as PlantVillage~\cite{hughes2015open}, where CNN and YOLO-based systems routinely achieve high accuracy~\cite{mohanty2016using, ferentinos2018deep}. However, real-world field deployment exposes two structural failures of these closed-set models: (1) a domain gap between laboratory-style training images and unconstrained field photographs, and (2) the inability to reject diseases or pests outside the training vocabulary. When a closed-set detector encounters an unknown class, its softmax output still assigns a confident label --- a property Dhamija et al.~\cite{dhamija2020overlooked} term closed-set overconfidence.

This paper investigates whether separating localization from open-set classification can address both failure modes. We use Grounding DINO~\cite{liu2023grounding} for language-guided region proposals and a fine-tuned CLIP model~\cite{radford2021learning} for per-crop verification with confidence thresholding. To motivate the design, we first present an empirical failure analysis of YOLOv11~\cite{jocher2024ultralytics} trained on PlantVillage-YOLO, identifying three distinct failure patterns that tiling strategies (SAHI~\cite{wang2022sahi}) do not resolve.

The main contributions of this paper are:
\begin{enumerate}
    \item An empirical failure analysis of YOLOv11 on real-world plant images, identifying scale mismatch, closed-set overconfidence, and tiling artifacts as three distinct failure modes.
    \item A multimodal open-set pipeline that couples Grounding DINO with fine-tuned CLIP and confidence thresholding, evaluated against a ResNet-50 closed-set baseline.
    \item An open-set evaluation protocol based on Unknown F1, which balances false rejection and missed unknowns under a threshold sweep.
    \item A few-shot extension for registering new disease prototypes at runtime using frozen CLIP embeddings and cosine similarity.
\end{enumerate}

\section{Related Work}

\subsection{Supervised Plant Disease Detection}
CNN-based approaches set the early standard for plant disease classification. Mohanty et al.~\cite{mohanty2016using} reported near-human accuracy on PlantVillage using AlexNet and GoogLeNet, and Ferentinos~\cite{ferentinos2018deep} extended this to deeper architectures. Object detection followed: single-stage YOLO detectors became widely adopted for field scouting due to their real-time throughput~\cite{li2021comprehensive}. However, Barbedo~\cite{barbedo2018impact} showed that these models often rely on shallow cues such as leaf shape and background color rather than the pathology itself --- performance degrades sharply when images move from controlled lab settings to open fields. Critically, all of these models operate under a closed-set assumption: the label space is fixed at training time, and any out-of-distribution input is silently misclassified.

A deeper issue underlies this: the lack of large-scale, real-world field datasets. PlantVillage~\cite{hughes2015open}, with over 54,000 images across 38 plant-disease categories, remains the dominant benchmark in plant pathology. However, every image in PlantVillage is a close-up photograph of a \emph{single isolated leaf}, taken under controlled lighting against a uniform background. Real field conditions --- full plants with multiple leaves at varying scales, partial occlusion, complex backgrounds, and inconsistent lighting --- are absent from the dataset entirely. Smaller datasets such as PlantDoc~\cite{singha2020plantdoc} attempt to close this gap with web-scraped field images, but they remain limited in scale and class coverage. As a result, the overwhelming majority of published models are trained and validated on close-up imagery and inherit a strong implicit prior about object scale and scene composition. The scale mismatch failure described in Section~\ref{sec:motivation} is therefore not a model-specific anomaly --- it is a structural consequence of training on the available data.

\subsection{Open-Set Recognition}
The open-set recognition problem~\cite{dhamija2020overlooked} requires a model to identify when an input falls outside its training vocabulary rather than forcing a classification. In the vision domain, standard approaches include OpenMax, which fits class-conditional extreme value distributions on activation vectors, and threshold-based rejection on maximum softmax probability (MSP). Scheirer et al. originally formalized the open-set constraint as learning a function that maps inputs to a recognition class or an explicit ``unknown'' label. Our work adapts MSP thresholding to the CLIP embedding space, exploiting the observation that CLIP confidence scores carry more semantic meaning than those of closed-set classifiers.

\subsection{Vision-Language Models in Agriculture}
CLIP~\cite{radford2021learning} learns joint image-text embeddings via contrastive pre-training on internet-scale data, giving it strong zero-shot transfer. Fine-tuning on domain-specific data preserves this open-set capability while improving accuracy on known classes~\cite{clip_plant_2024}. Grounding DINO~\cite{liu2023grounding} extends the DINO detector with cross-modal attention, enabling object localization driven by free-form text prompts rather than fixed class indices. Recent studies on agricultural datasets report that Grounding DINO achieves higher mAP than fine-tuned YOLO under few-shot conditions, and it has been used as an auto-labeling teacher to reduce annotation costs. Together, these properties make it well suited for detecting leaves and pests without disease-specific bounding box annotations.

\subsection{Few-Shot and Runtime Adaptation}
Registering new categories at test time without retraining is an active research area. Prototype-based methods~\cite{snell2017prototypical} compute class centroids in an embedding space and classify by nearest centroid. Vision-Language frameworks such as VLCD apply CLIP features as a ``cache model'' to blend prior knowledge with few-shot task examples. Our few-shot engine follows the same principle: frozen CLIP embeddings are averaged into per-disease centroids, and cosine similarity against these centroids classifies regions that fall below the confidence threshold. This allows new disease types to be registered from a small number of images at runtime.

\section{Motivation: Failure of Closed-Set Detectors}
\label{sec:motivation}

To motivate the design of our pipeline, we trained YOLOv11~\cite{jocher2024ultralytics} on a bounding-box-annotated version of PlantVillage~\cite{hughes2015open} (sourced from Kaggle~\cite{plantvillage_yolo_kaggle}) and tested it on real field photographs of full plants with multiple diseased leaves. Three consistent failure patterns emerged.

\subsection{Finding 1: Scale Mismatch}

The PlantVillage dataset consists of close-up, single-leaf photographs in which the leaf fills nearly the entire frame. Trained on this data, YOLOv11 learns a strong spatial prior: one object, centered, occupying most of the image. When applied to a real full-plant photograph containing multiple leaves at varying distances and scales, the model returned a \emph{single} bounding box spanning the entire image area. It had no mechanism to decompose a multi-object scene --- it defaulted to the dominant prior learned from close-up training images. This is a classic instance of domain shift driven by object scale distribution mismatch between training and deployment.

\subsection{Finding 2: Closed-Set Overconfidence}

The scale-mismatch failure is compounded by the model's inability to express uncertainty. YOLOv11 uses a softmax classification head, which forces all probability mass across the fixed set of known training classes regardless of input quality. On the full-plant test images, the model produced a confident disease label even when the bounding box covered the entire scene --- a combination of location error and forced misclassification. There is no threshold or score below which the model will decline to classify; it always returns a winner. This means errors are silent and appear reliable, which is dangerous in a deployment context where an unknown or ambiguous case should trigger human review.

\subsection{Finding 3: SAHI Limitations on Irregular Objects}

Slicing Aided Hyper Inference (SAHI)~\cite{wang2022sahi} addresses scale mismatch by tiling the input image into a grid of overlapping fixed-size patches and running the detector on each patch independently. In principle, smaller tiles better match the close-up scale of the training data. In practice, we found two problems on real plant photography. First, many tiles land entirely on background --- soil, open sky, or homogeneous foliage --- and the model still predicts a leaf bounding box with nonzero confidence on these empty patches. Second, tiles that do contain a leaf frequently cut across the natural leaf boundary, producing detections that do not correspond to a meaningful region. Unlike standard SAHI benchmarks where objects are small but regularly shaped (e.g., pedestrians, vehicles), agricultural leaves vary widely in size, shape, orientation, and overlap. No fixed tile size is appropriate across this distribution, and the closed-set overconfidence problem (Finding 2) persists inside each tile regardless.

These three findings directly motivate the Grounding DINO + CLIP design: language-guided detection removes the scale prior, and confidence thresholding provides an explicit rejection mechanism for out-of-distribution inputs.

\section{Proposed System}

The pipeline follows an \textit{Observe~$\rightarrow$~Verify~$\rightarrow$~Reason~$\rightarrow$~Act} loop. Figure~\ref{fig:pipeline} gives an overview. Localization and identification are kept strictly separate: Grounding DINO finds regions of interest in the full image, CLIP classifies each cropped region, and an LLM agent synthesises the per-region findings into a final diagnosis.

\subsection{Stage 1: Open-Vocabulary Detection (Grounding DINO)}
\label{subsec:gdino}

We use Grounding DINO (Swin-B backbone)~\cite{liu2023grounding} with a fixed text prompt ``\texttt{leaf . bug . worm}'' to localize all agriculturally relevant regions simultaneously. This prompt-driven approach requires no disease-specific bounding box annotations and generalises to unseen crop types without retraining.

Two detection thresholds control the recall/precision trade-off:
\begin{itemize}
    \item \textbf{Box threshold} $\tau_{\text{box}} = 0.30$: a region is retained only if the model is at least 30\% confident an object is present.
    \item \textbf{Text threshold} $\tau_{\text{text}} = 0.25$: the detected visual feature must align with the semantic token in the prompt.
\end{itemize}

Detected bounding boxes are cropped and normalized before being passed to the verification stage. The choice of Grounding DINO here is an implementation decision --- the architecture is modular, and any open-vocabulary detector (e.g., YOLO-World~\cite{wang2024yoloworld}) could serve the same localization role. The open-set rejection capability of the pipeline derives from the CLIP verification stage, not from the specific detector used.

\subsection{Stage 2: Open-Set Verification (CLIP)}
\label{subsec:clip}

Each crop is classified by one of two CLIP variants, selected based on the detection label:

\begin{itemize}
    \item \textbf{Leaf regions}: a fine-tuned CLIP model (\texttt{Keetawan/clip-vit-large-patch14-plant-disease-finetuned}~\cite{keetawan2024plantdisease}) is used for disease identification. Fine-tuning preserves zero-shot generalization while improving accuracy on the 10 target Tomato pathologies.
    \item \textbf{Pest regions}: the base \texttt{openai/clip-vit-base-patch16} is used. The general model avoids the leaf-feature bias introduced by domain fine-tuning and identifies insects (e.g., aphids, grasshoppers) from their natural language descriptions.
\end{itemize}

Classification uses the maximum softmax probability (MSP) over all label embeddings. For an image embedding $I$ and label embeddings $\{T_k\}_{k=1}^{K}$, the confidence score is:
\begin{equation}
    p_k = \frac{\exp(\text{sim}(I, T_k)/\tau)}{\sum_{j=1}^{K} \exp(\text{sim}(I, T_j)/\tau)}
    \label{eq:clip_softmax}
\end{equation}
where $\tau$ is CLIP's learned temperature and $\text{sim}(\cdot,\cdot)$ is cosine similarity. A prediction is accepted only if $\max_k p_k \geq \delta$; otherwise the region is flagged as \texttt{Unknown} and routed to the few-shot engine.

\subsection{Stage 2b: Few-Shot Engine for Unknown Classes}
\label{subsec:fewshot}

Rather than discarding unrecognized regions, the system attempts to match them against user-registered disease prototypes. The few-shot engine stores per-class centroids computed from the frozen CLIP image encoder:
\begin{equation}
    \mu_c = \frac{1}{|S_c|} \sum_{x \in S_c} \phi(x)
\end{equation}
where $S_c$ is the support set for class $c$ and $\phi(\cdot)$ is the CLIP image encoder. At inference, a region embedding $\phi(x_{\text{test}})$ is assigned to class $c^*$ if:
\begin{equation}
    c^* = \arg\max_c \; \text{sim}(\phi(x_{\text{test}}), \mu_c), \quad \text{if} \; \max_c \; \text{sim} \geq \theta_{\text{fs}}
\end{equation}
with similarity threshold $\theta_{\text{fs}} = 0.82$. If no prototype exceeds this threshold, the region retains the \texttt{Unknown disease} label. Prototypes are registered at runtime from as few as one image per class and persisted to disk, requiring no model retraining.

\subsection{Stage 3: Reasoning and Diagnosis (LLM Agent)}
\label{subsec:agent}

The agent (Gemini 2.5 Flash, orchestrated via Pydantic AI) receives a structured summary of all verified detections --- disease labels, confidence scores, bounding box counts, and the infection ratio per class. It invokes a RAG tool (\texttt{consult\_ipm\_manual}) that queries a ChromaDB vector store of agricultural manuals, filtered by the detected crop species, to retrieve treatment protocols. The agent's output is constrained to a typed \texttt{DiagnosisResult} schema enforced by Pydantic, which includes three self-correction validators that trigger a model retry if the diagnosis is internally inconsistent (e.g., high severity with no recommended treatment).

\section{Experiments}
\subsection{Experimental Setup}
We evaluate the system on 10 known Tomato disease classes (100 samples each) and unknown classes such as Potato Late blight. We compare our pipeline against a ResNet-50 baseline.

\subsection{Results}
Table~\ref{tab:results} summarizes the performance at a confidence threshold of 0.35. Table~\ref{tab:sweep} demonstrates the effect of varying the threshold.

\begin{table}[h]
\centering
\caption{Performance Comparison at Threshold = 0.35}
\label{tab:results}
\begin{tabular}{@{}lcc@{}}
\toprule
Metric & CLIP (Open-Set) & ResNet-50 (Closed-Set) \\ \midrule
Known Accuracy & 80.0\% & Baseline \\
Unknown Rejection Rate & 40.0\% & 50.0\% \\
False Rejection Rate & 13.3\% & 63.3\% \\
Unknown F1 & 0.444 & 0.294 \\ \bottomrule
\end{tabular}
\end{table}

\begin{table}[h]
\centering
\caption{CLIP Threshold Sweep}
\label{tab:sweep}
\begin{tabular}{@{}ccccc@{}}
\toprule
Threshold & Known Acc & Unk Rejection & False Rejection & Unknown F1 \\ \midrule
0.4 & 80.0\% & 100.0\% & 10.0\% & 0.375 \\
0.5 & 80.0\% & 40.0\% & 13.3\% & 0.444 \\
\textbf{0.6} & \textbf{80.0\%} & \textbf{50.0\%} & \textbf{13.3\%} & \textbf{0.526} \\
0.7 & 73.3\% & 50.0\% & 20.0\% & 0.476 \\ \bottomrule
\end{tabular}
\end{table}

\section{Discussion}
Selecting the optimal confidence threshold involves a trade-off between false rejections (which degrade user experience) and missing unknown diseases. At a threshold of 0.6, we achieved the best Unknown F1 score.

\section{Conclusion and Future Work}
We introduced a multimodal open-set pipeline using Grounding DINO and CLIP to address the limitations of closed-set detectors in agricultural AI. Future work will involve expanding the dataset size, resolving baseline evaluation anomalies, and evaluating the integration of Retrieval-Augmented Generation (RAG)~\cite{lewis2020retrieval} for providing treatment recommendations.

\bibliographystyle{splncs04}
\bibliography{references}

\end{document}
